%=====================================================================
% Cover letter. Compile with: python src/build_journal.py
%=====================================================================
\documentclass[11pt]{article}
\usepackage[a4paper,margin=2.5cm]{geometry}
\usepackage{amsmath,amssymb}
\usepackage{hyperref}
\pagestyle{empty}

\begin{document}

\begin{flushright}
Jingwen Zhang, Fei Ma\\
School of Mathematics and Physics\\
Xi'an Jiaotong-Liverpool University\\
Suzhou, China
\end{flushright}

\vspace{1em}
\noindent Dear Editors of the \emph{Journal of Computational Physics},

\vspace{0.6em}
\noindent We submit for your consideration the manuscript \emph{``Reconstruction Order, Not Solver
Family, Governs Control Accuracy: A Comparison of Neural and Classical HJB Solvers Against an Exact
Reference''}.

\noindent The paper measures something that is usually asserted. Neural solvers for
Hamilton--Jacobi--Bellman equations are proposed on the strength of the curse of dimensionality, and
the comparison literature scores them against other solvers --- which makes the reference exact by
construction and its reported error zero. We score four families against a closed-form Riccati
reference on a multi-asset optimal-execution HJB, and we report the \emph{control} as well as the
value. Three results follow, and each is stated as a rate rather than an observation:

\begin{itemize}
\item \textbf{Two propositions.} A degree-$k$ reconstruction delivers gradient order $k$, but a
symmetric-difference probe at step $\varepsilon$ adds an independent $O(\varepsilon^2)$, so a probe
that shrinks only like $h$ caps a cubic reconstruction at order two. And for this problem the
objective lost is \emph{exactly} the $R$-weighted squared control deviation accumulated along the path
taken --- an identity, not the bound the envelope theorem would give. Together they explain why the
control error is 3.2 to 26 times the value error, and they reduce to a prescription: raise the
reconstruction order, and scale the probe with it.
\item \textbf{A boundary, measured rather than assumed.} Where a grid scheme runs it is faster and
more accurate; we also run a compiled kernel of the same scheme, which changes the comparison at the
top dimension, and we run the classical mesh-free family the earlier version of this work had only
cited --- where it turns out to be \emph{more} accurate than the neural solver at every dimension
tested. A conclusion changed as a result, and the paper says so.
\item \textbf{What it is worth.} On the paper's own real-market calibration, raising the
reconstruction from multilinear to cubic is worth 2.7 to 8.4 basis points of notional across nine
instruments, with the largest effect on the most liquid one.
\end{itemize}

\noindent The paper is a comparison and a measurement, and we have tried to make it useful in the way
your readership would want: every table is generated from committed result files, the verification
suite runs on a clean checkout, the scheme's own solver is the one used for the compiled timing so
that implementation and algorithm are separated, and the limitations are stated where they bind
rather than collected in a closing paragraph. Three of them are worth naming here. The grid sweep
stops at $p=5$; the neural results use one architecture family at stated budgets; and the
non-polynomial benchmark we construct to answer the ``your solution is quadratic'' objection uses a
manufactured source, which produces an exact solution but artificial economics.

\noindent We confirm that this manuscript is original, is not under consideration elsewhere, and that
all authors have approved the submission. Declarations of competing interest, the CRediT statement
and the data availability statement accompany the submission. Reproducibility is the point of the
paper, so we have made every number traceable to a released file, and the repository is archived.

\vspace{1em}
\noindent We would be glad to suggest referees if that is helpful, and we thank you for considering
the work.

\vspace{1.4em}
\noindent Sincerely,\\
Jingwen Zhang and Fei Ma

\end{document}
